\documentclass[../main.tex]{subfiles}
\begin{document}

\section{Problem Statement}
\label{sec:dvd}
Legal consultation and legal judgement prediction have become important research problems in the legal domain. In practice, understanding a legal case requires not only reading the facts of the case, but also identifying relevant legal articles, interpreting the relationship between legal provisions and case circumstances, and considering factors that may affect the final judgement. For ordinary citizens, this process is difficult because legal documents are usually written in formal language and require professional knowledge to navigate and interpret correctly. 

In recent years, artificial intelligence, especially large language models, has shown strong potential in supporting legal tasks such as legal question answering, case analysis, and judgement prediction. LLMs are capable of understanding natural language queries, performing high-level reasoning, and identifying fine-grained details within complex textual contexts, which makes them suitable for building legal consultation systems. However, using LLMs directly for legal judgement prediction still has several limitations. Since LLMs mainly rely on their pre-trained knowledge, they may provide outdated or incomplete legal information. LLMs are also prone to hallucination without explicit instruction and detailed law article text. This problem is more serious in the legal domain, where accuracy, traceability, and consistency with official legal documents are essential.

Compared with English legal systems, Vietnamese legal resources are less available in structured and machine-readable formats. Court judgements are often published as PDF files, which makes automatic extraction and analysis more difficult. Vietnamese legal judgement prediction also requires the system to connect case facts with relevant law articles, criminal charges, penalty ranges, and mitigating or aggravating circumstances. In many cases, legal violations may appear highly similar and are separated only by subtle differences in violation intention or factual circumstances. Therefore, a simple text generation model is usually not sufficient to perform this task reliably.

This thesis focuses on addressing this problem by developing a retrieval-augmented generation system combined with a multi-stage retrieval-grounded reasoning pipeline for Vietnamese criminal judgement prediction. The proposed system aims to support the prediction process by retrieving relevant legal articles and similar past cases, then guiding the LLM to reason over the case facts in a more structured and traceable manner. By combining statutory law, case law retrieval, and step-by-step reasoning, the system is expected to reduce hallucination and improve the reliability of predictions for criminal charges and penalty lengths. This approach also provides a foundation for building more practical legal AI systems for the Vietnamese legal domain.

\section{Objectives and Scope } 
\label{sec:giaiphap}
AI systems is capable of supporting legal tasks such as charge prediction, legal article recommendation, and penalty estimation. In particular, LLM-based methods are promising because they can understand natural language queries, perform high-level reasoning, and identify fine-grained details within complex legal texts.However, existing approaches still face several limitations. First, many legal judgement prediction systems are developed for languages and legal systems with large-scale structured datasets, such as English or Chinese, while Vietnamese legal judgement prediction remains less explored. Second, direct use of LLMs may lead to hallucination if the model is not grounded in detailed legal articles and reliable case law data. Third, criminal offences can be highly similar and are sometimes separated only by subtle differences in legal elements or factual circumstances. Therefore, a simple text generation model is usually not sufficient to perform this task reliably.

Based on these limitations, this thesis aims to develop a legal judgement prediction system for Vietnamese criminal law. The system focuses on supporting the prediction of criminal charges and penalty lengths from structured case information. Instead of treating judgement prediction as a pure text generation task, the thesis investigates how external legal knowledge and structured reasoning can improve the reliability of LLM-based prediction.

The main objectives of this thesis are:

\begin{enumerate}
    \item To construct a structured dataset of Vietnamese criminal judgements for retrieval and evaluation.
    \item To develop a system that supports legal article retrieval, similar case retrieval, and LLM-based judgement prediction.
    \item To predict criminal charges and penalty lengths from defendant information and case summaries.
    \item To evaluate the proposed system against LLM-only methods using ground-truth judgements in the test dataset.
\end{enumerate}

The scope of this thesis is limited to Vietnamese criminal law. The input of the system mainly consists of defendant information and case summaries extracted from criminal judgement documents. The expected outputs are predicted criminal charges and penalty lengths. The thesis does not cover other legal domains such as civil, administrative, commercial, or labor law.

This thesis is developed as a research prototype and does not aim to replace professional legal advice or the decision-making role of judges. Its purpose is to examine whether retrieval and structured reasoning can improve the performance of LLMs in Vietnamese legal judgement prediction. The system performance is also limited by the quality, coverage, and structure of the collected legal data.

\section{Proposed Solution}
To address the problems discussed in the previous section, this thesis proposes a retrieval-augmented generation system combined with a multi-stage retrieval-grounded reasoning pipeline for Vietnamese criminal judgement prediction. The main goal of the proposed solution is to support the prediction of criminal charges and penalty lengths by combining three sources of information: the facts of the input case, relevant legal articles, and similar past court judgements. Instead of relying only on the pre-trained knowledge of a large language model, the system retrieves legal knowledge from external databases and guides the model through a structured reasoning process.

First, to address the lack of structured Vietnamese legal case data, this thesis builds a criminal case law database from raw judgement PDF files collected from the Supreme People's Court of Vietnam. Since these documents are published in PDF format, the system uses Optical Character Recognition (OCR) to extract text from the files. The extracted text is then cleaned and organized into structured JSON records. Important fields such as defendant information, case summary, court reasoning, applied legal articles, criminal charges, and penalties are extracted and stored for later retrieval and evaluation. This database forms the foundation for both training data analysis and similar case retrieval.

Second, to reduce the risk of hallucination and incomplete legal reasoning, the system integrates a retrieval-augmented generation approach. Relevant legal articles and past criminal cases are indexed in a vector database using an embedding model. When a new case is provided, the system retrieves legal articles that are related to the case facts, as well as past cases with similar circumstances. These retrieved documents provide the LLM with concrete legal evidence and examples, allowing it to generate predictions based on available legal materials instead of relying only on its internal knowledge.

Third, to handle the complexity of legal reasoning, this thesis designs a multi-stage retrieval-grounded reasoning pipeline for Vietnamese criminal judgement prediction. The pipeline is inspired by the Reasoning-and-Act idea of separating reasoning from external actions, but the sequence of actions is predefined by the system rather than dynamically selected by the LLM. In the implemented system, the pipeline is organized as three structured stages rather than a single direct generation prompt.

In the first stage, the LLM receives the defendant information and synthetic case summary, then extracts structured facts from the case. It identifies the defendants, relevant conduct, property or harm information, prior convictions, mitigating signals, aggravating signals, stated facts, inferred facts, and missing facts. Based on these facts, the model proposes several plausible candidate offences under the Vietnamese Criminal Code. This stage is used to prevent the system from committing immediately to one charge without first considering alternative legal possibilities.

In the second stage, the system acts on the candidate offences by retrieving the corresponding legal articles from the statutory law database. In addition to offence-specific articles, the system always checks a fixed set of supporting criminal-law articles, including articles on imprisonment, sentencing basis, mitigating circumstances, aggravating circumstances, recidivism, multiple offences, attempted offences, accomplice liability, suspended sentence, and property handling. The LLM then performs legal analysis over the extracted facts, the candidate offences, and the retrieved law text. It selects the most likely offence, identifies the sentencing bracket, rejects or downgrades unsuitable candidates, and classifies supporting articles as applicable, fact-dependent, not applicable, or not retrieved.

In the third stage, the system retrieves similar past cases that match the selected offence article. These cases are used only for analogy and sentencing calibration, while statutory law remains the controlling source for charge and sentencing-frame selection. The final LLM call receives the original case fields, the legal analysis, retrieved offence articles, retrieved supporting articles, and summaries of similar cases. It then produces the final structured prediction, including the predicted charge, applied legal clauses, imprisonment penalty, additional fine, civil liability, and handling of physical evidence when available.

Fourth, to distinguish between legally similar violations, the system explicitly compares the factual elements of the case with the legal elements of candidate offences. Many criminal violations may appear similar but differ in subtle details such as intent, damage level, method of committing the offence, or the role of each defendant. Therefore, the proposed system requires the LLM to justify why a specific charge is selected and why other candidate charges are less suitable. This comparison-based reasoning is expected to improve the reliability of charge prediction.

Finally, the system evaluates its predictions by comparing them with the ground-truth judgements in the test dataset. The evaluation focuses mainly on legal article prediction and penalty length prediction. Legal article prediction is evaluated using precision, recall, and F1 score after normalizing predicted and ground-truth legal clauses to article level. Penalty prediction is evaluated by converting imprisonment sentences into months and computing the root mean squared error. The proposed pipeline is compared with simpler LLM-only and retrieval-based baselines in order to measure whether structured reasoning and grounded retrieval improve judgement prediction.

In summary, the proposed solution consists of a complete pipeline for Vietnamese criminal legal judgement prediction. It begins with collecting and structuring raw legal judgement PDFs, continues with building a searchable legal and case law database, and ends with an LLM-based reasoning system that retrieves, analyzes, and predicts legal outcomes. This approach is designed to improve factual grounding, reduce hallucination, and provide a more transparent reasoning process for legal judgement prediction in the Vietnamese criminal law domain.

\section{Contributions}
\begin{enumerate}
\item Implement a complete pipeline for collecting Vietnamese criminal judgement PDF files and converting them into structured JSON records. The structured records include defendant information, case facts, court reasoning, legal articles, criminal charges, penalties, and other judgement-related fields.
\item Construct a retrieval database for Vietnamese criminal law and past criminal cases. The database supports both statutory-law lookup and similar-case retrieval, which are used to ground LLM-based judgement prediction.
\item Propose a multi-stage retrieval-grounded reasoning pipeline for Vietnamese criminal judgement prediction. The pipeline decomposes prediction into fact extraction and candidate generation, statutory legal analysis, similar-case retrieval, and final structured verdict generation.
\end{enumerate}

\end{document}
